\documentclass[12pt]{article}

% Import preambles and macros for homework
\input{../../../../preambles/hw-preamble.tex}
\input{../../../../preambles/homework-macros.tex}
\input{../../../../preambles/homework-title.tex}
\input{../../../../preambles/homework-config.tex}

\renewcommand{\author}{Deepak Jassal}
\renewcommand{\authorlast}{Jassal}
\renewcommand{\coursename}{Advanced Linear Algebra}
\renewcommand{\coursecode}{MATH 326}
\renewcommand{\assignment}{Assignment 5}
\renewcommand{\instructor}{Dr. Edward Dobrowolski}
\renewcommand{\duedate}{March 6th\textsuperscript{th}, 2026}
\usepackage{tensor}

\begin{document}
\begin{titlepage}
	\centering
	\vspace*{2.0cm}	
	\pgfornament{84}\\
	{\LARGE \textsc{\coursename}\par}
	\vspace{0.5cm}
	{\large\coursecode\par}
    \vspace{0.5cm}
    {\large\instructor\par}
	\vspace{1.5cm}
	{\huge\bfseries\assignment\par}
	\vspace{1cm}  
	{\LARGE\itshape\author\par}
    \vspace{2cm}
	{\large\bfseries Due Date:\par}
	\vspace{0.5cm}
	{\Large \duedate}\\
	\pgfornament{84}
\end{titlepage}
\stepcounter{section}
\section*{Problem 1}
Let $\mathcal{V}$ be a real inner product and let $x, y \in \mathcal{V}$. Show that $\inner{x}{y}= 0$ if and only if $\norm{x+y}=\norm{x-y}$.
\begin{proof}
    $(\Rightarrow)$ Assume that $\inner{x}{y}=0$, then
    \begin{align*}
        \norm{x+y}^2&=\inner{x+y}{x+y}\\
        &=\inner{x}{x}+\inner{x}{y}+\inner{y}{x}+\inner{y}{y}\\
        &=\inner{x}{x}+2\inner{x}{y}+\inner{y}{y}\\
        &=\inner{x}{x}-2\inner{x}{y}+\inner{y}{y}
        \intertext{Since $\inner{x}{y}=-\inner{x}{y}=0.$}
        &=\norm{x-y}^2\\
        \norm{x+y}&=\norm{x-y}.
    \end{align*}
    $(\Leftarrow)$ Assume that $\norm{x+y}=\norm{x-y}$, then
    \begin{align*}
        \norm{x+y}^2&=\norm{x-y}^2\\
        \inner{x}{x}+2\inner{x}{y}+\inner{y}{y}&=\inner{x}{x}-2\inner{x}{y}+\inner{y}{y}\\
        2\inner{x}{y}&=-2\inner{x}{y}\\
        \inner{x}{y}&=0.\qedhere
    \end{align*}
\end{proof}
\stepcounter{section}
\section*{Problem 2}
Let $\mathcal{V}$ be an $\mathbb{F}$-product space and let $u, v \in \mathcal{V}$. Show that
\[
    |\inner{u}{v}|\leq \norm{u}^2+\norm{v}^2.
\]
When is this inequality an equality?\\
\textit{Solution.} Let $a=\norm{u}^2\geq0$ and $b=\norm{v}^2\geq0$. By the Cauchy-Schwarz-Bunyakovsky inequality we have
\begin{align*}
    |\inner{u}{v}|&\leq\norm{u}\norm{v}\\
    |\inner{u}{v}|&\leq\sqrt{a}\sqrt{b}.
\end{align*}
By the AM-GM inequality we have
\[
    |\inner{u}{v}|\leq\sqrt{a}\sqrt{b}=\sqrt{ab}\leq\frac{a+b}{2}\leq a+b=\norm{u}^2+\norm{v}^2.
\]
Thus we have shown the first part. To see when this is an equality note that
\[
    |\inner{u}{v}|=a+b\Rightarrow|\inner{u}{v}|=\frac{a+b}{2}\Rightarrow\frac{a+b}{2}=a+b,
\]
from here see that 
\begin{align*}
    \frac{a+b}{2}&=a+b\\
    a+b&=0,
\end{align*}
since $a,b\geq0$, we have equality whenever
\[
    \norm{u}^2=\norm{v}^2=0.
\]
\stepcounter{section}
\section*{Problem 3}
Let $\mathcal{V}$ be a complex vector space and let $u, v \in \mathcal{V}$. Show that
\[
    \inner{u}{v}=\frac{1}{2\pi}\int_{-\pi}^{\pi}e^{i\theta}\norm{u+e^{i\theta}v}^2\,d\theta.
\]
\begin{align*}
    \inner{u}{v}&=\frac{1}{2\pi}\int_{-\pi}^{\pi}e^{i\theta}\norm{u+e^{i\theta}v}^2\,d\theta\\
    &=\frac{1}{2\pi}\int_{-\pi}^{\pi}e^{i\theta}\inner{u+e^{i\theta}v}{u+e^{i\theta}v}\,d\theta\\
    &=\frac{1}{2\pi}\int_{-\pi}^{\pi}e^{i\theta}\left(\inner{u}{u}+\inner{e^{i\theta}v}{u}+\inner{u}{e^{i\theta}v}+\inner{e^{i\theta}v}{e^{i\theta}v}\right)\,d\theta\\
    &=\frac{1}{2\pi}\int_{-\pi}^{\pi}e^{i\theta}\left(\norm{u}^2+e^{i\theta}\inner{v}{u}+e^{-i\theta}\inner{u}{v}+e^{i\theta}e^{-i\theta}\inner{v}{v}\right)\,d\theta\\
    &=\frac{1}{2\pi}\int_{-\pi}^{\pi}e^{i\theta}\left(\norm{u}^2+e^{i\theta}\inner{v}{u}+e^{-i\theta}\inner{u}{v}+\norm{v}^2\right)\,d\theta\\
    &=\frac{\norm{u}^2+\norm{v}^2}{2\pi}\int_{-\pi}^{\pi}e^{i\theta}\,d\theta+\frac{\inner{v}{u}}{2\pi}\int_{-\pi}^{\pi}e^{2i\theta}\,d\theta+\frac{\inner{u}{v}}{2\pi}\int_{-\pi}^{\pi}\,d\theta\\
    &=0+0+\frac{2\pi\inner{u}{v}}{2\pi}\\
    &=\inner{u}{v}.
\end{align*}
\newpage
\stepcounter{section}
\section*{Problem 4}
Let $\mathcal{V}$ be the real product space $C_\R[0, 1]$, let $f \in \mathcal{V}$, and suppose that $f (t) > 0$ for all $t \in [0, 1]$. Use the Cauchy-Schwarz-Bunyakovsky inequality to deduce that
\[
    \frac{1}{\int_{0}^{1}f(t)\,dt}\leq\int_{0}^{1}\frac{1}{f(t)}\,dt.
\]
\begin{proof}
    Let
    \[
        \inner{\phi}{\psi}=\int_{0}^{1}\phi(t)\psi(t)\,dt
    \]
    and
    \[
        \norm{\phi}=\left(\int_{0}^{1}\phi(t)^2\,dt\right)^{\frac{1}{2}}
    \]
    for $\phi,\psi\in C_\R[0, 1]$. Let $g(t)=\sqrt{f(t)}$ then by the Cauchy-Schwarz-Bunyakovsky inequality we have
    \begin{align*}
        |\inner{g}{g^{-1}}|&\leq\norm{g}\norm{g^{-1}}\\
        \left|\int_{0}^{1}g(t)g^{-1}(t)\,dt\right|&\leq\left(\int_{0}^{1}g^2(t)\,dt\right)^{\frac{1}{2}}\left(\int_{0}^{1}g^{-2}(t)\,dt\right)^{\frac{1}{2}}\\
        \left|\int_{0}^{1}1\,dt\right| &\leq\left(\int_{0}^{1}f(t)\,dt\right)^{\frac{1}{2}}\left(\int_{0}^{1}\frac{1}{f(t)}\,dt\right)^{\frac{1}{2}}\\
        1&\leq\int_{0}^{1}f(t)\,dt\int_{0}^{1}\frac{1}{f(t)}\,dt\\
        \frac{1}{\int_{0}^{1}f(t)\,dt}&\leq\int_{0}^{1}\frac{1}{f(t)}\,dt.\qedhere
    \end{align*} 
\end{proof}
\stepcounter{section}
\section*{Problem 5}
Let $x^T = [x_1, x_2, x_3]$ and $y^T = [y_1, y_2, y_3]$ be vectors in $\R^3$. Further let $f (x, y) = 12x_1y_1 +4x_1y_2 +3x_1y_3 +4x_2y_1 +2x_2y_2 -x_2y_3 +3x_3y_1 -x_3y_2 +7x_3y_3$.
\begin{enumerate}[label=(\alph*)]
    \item Represent $f$ in the form
    \[
        f(x,y)=x^TAy,
    \]
    where $A \in M_3$, and prove that $f (x, y)$ is an inner product on $\R^3$.
    \begin{proof}
        \[
            A=
            \begin{bmatrix}
                12&4&3\\
                4&2&-1\\
                3&-1&7
            \end{bmatrix}.
        \]
        Note that $A$ is a symmetric matrix, therefore all of its eigenvalues are real and positive. This gives that $f(x,y)$ is an inner product. Also each minor determinant of $A$ is positive, which by Sylvesters criterion gives that $f(x,y)$ is an inner product. 
    \end{proof}
    \item Find the formula for the norm of $x$ in the norm generated by $f$.\\
    \textit{Solution.} Let $x=[x_1,x_2,x_3]^T$. We know that 
    \[
        \norm{x}=\sqrt{\inner{x}{x}}=\sqrt{f(x,x)}.
    \]
    From here we have
    \begin{align*}
        f(x,x)&=12x_1x_1 +4x_1x_2 +3x_1x_3 +4x_2x_1 +2x_2x_2 -x_2x_3 +3x_3x_1 -x_3x_2 +7x_3x_3\\
        &=12x_1^2+4x_1x_23+3x_1x_3 +4x_2x_1 +2x_2^2 -x_2x_3 +3x_3x_1 -x_3x_2 +7x_3^2\\
        \norm{x}&=\sqrt{12x_1^2+4x_1x_23+3x_1x_3 +4x_2x_1 +2x_2^2 -x_2x_3 +3x_3x_1 -x_3x_2 +7x_3^2}.
    \end{align*}
\end{enumerate}
\stepcounter{section}
\section*{Problem 6}
Let $\mathcal{P}_{n-1}$ be the space of all real polynomials of degree not exceeding $n - 1$. If $p(x)=\sum_{k=0}^{n-1}p_kx^k$ and $q(x)=\sum_{k=0}^{n-1}q_kx^k$, define
\[
    \inner{p}{q}=\sum_{i,j=0}^{n-1}\frac{p_iq_j}{i+j+1}.
\]
\begin{enumerate}[label=(\alph*)]
    \item Show that $\inner{p}{q}$ is an inner product on $\mathcal{P}_{n-1}$.\\
    \textit{Hint.} We can consider $\mathcal{P}_{n-1}$ as a subspace of $C_R[0, 1]$. Hence $\int_{0}^{1}p(x)q(x)\,dx$ is also an inner product on $\mathcal{P}_{n-1}$. Calculate this inner product on a few polynomials $p$ and $q$ and compare the results with the formula above.
    \begin{proof}
        Since addition and multiplication is commutative in $\R$ we have for $p,q\in\mathcal{P}$.
        \[
            \inner{p}{q}=\sum_{i,j=0}^{n-1}\frac{p_iq_j}{i+j+1}=\sum_{i,j=0}^{n-1}\frac{q_jp_i}{j+i+1}=\inner{q}{p}.
        \]
        So $\inner{p}{q}$ is symmetric.\\
        Let $a,b\in\R$ and $p,q,r\in\mathcal{P}_{n-1}$. Then,
        \begin{align*}
            \inner{ap+br}{q}&=\sum_{i,j=0}^{n-1}\frac{(ap_i+br_i)q_j}{i+j+1}\\
            &=\sum_{i,j=0}^{n-1}a\frac{p_iq_j}{i+j+1}+b\frac{r_iq_j}{i+j+1}\\
            &=a\sum_{i,j=0}^{n-1}\frac{p_iq_j}{i+j+1}+b\sum_{i,j=0}^{n-1}\frac{r_iq_j}{i+j+1}\\
            &=a\inner{p}{q}+b\inner{r}{q}.
        \end{align*}
        So $\inner{p}{q}$ is linear in the first slot.\\
        Let
        \[
            \inner{p}{q}_{L^2}=\int_{0}^{1}p(x)q(x)\,dx.
        \]
        We know that
        \[
            p(x)q(x)=\sum_{j,k=0}^{n-1}p_iq_jx^{i+j}.
        \]
        Substituting this into the integral
        \begin{align*}
            \inner{p}{q}_{L^2}&=\int_{0}^{1}p(x)q(x)\,dx\\
            &=\int_{0}^{1}\sum_{j,k=0}^{n-1}p_iq_jx^{i+j}\,dx
            \intertext{Since the sum is finite}
            &=\sum_{j,k=0}^{n-1}p_iq_j\int_{0}^{1}x^{i+j}\,dx\\
            &=\sum_{i,j=0}^{n-1}\frac{p_iq_j}{i+j+1}\\
            &=\inner{p}{q}.
        \end{align*}
        Since $\inner{p}{q}_{L^2}$ is an inner product, we have that $\inner{p}{q}$ is positive definate and therefore is an inner product.
    \end{proof}
    \item Deduce that the matrix $A = [(i + j - 1)^{-1}] \in M_n$ is positive definite.\\
    \textit{Deduction.} Let 
    \[
        p=\begin{bmatrix}
            p_1\\\vdots\\p_n
        \end{bmatrix},\quad
        q=\begin{bmatrix}
            q_1\\\vdots\\q_n
        \end{bmatrix}.
    \]
    \[
        pAq^T=\sum_{i,j=1}^{n}\frac{p_iq_j}{i+j-1}
    \]
    re-indexing the sum $i=i+1$ and $j=j+1$ we get
    \[
        pAq^T=\sum_{i,j=0}^{n-1}\frac{p_{i+1}q_{j+1}}{i+j1}=\inner{p}{q}.
    \]
    Since in part (a) we showed that $\inner{p}{q}$ is an inner product, that gives that $A$ is a positive definate matrix.
\end{enumerate}

\end{document}